\chapter{病句辨析与修改}

\begin{introduction}
  \item 搭配不当
  \item 用词不当
  \item 语序不当
  \item 成分残缺
  \item 成分赘余
  \item 结构混乱
  \item 不合逻辑
  \item 表意不明
\end{introduction}


\section{主要类型知识归纳}

\subsection{搭配不当}

\subsubsection{主谓搭配不当}

例：他一进教室，同学们的眼睛都集中到他的身上。

解析：句中主语“眼睛”与谓语“集中”搭配不当。“眼睛”不能“集中”，应改为“目光”。

\subsubsection{动宾搭配不当}

例：英雄们把红旗和胜利插上了敌人的阵地。

解析：句中动宾搭配不当，“红旗”可以插，“胜利”怎么插？应改为“英雄们把胜利的红旗插上了敌人的阵地”。

\subsubsection{主宾搭配不当}

例：冬天的上海是美丽的季节。

解析：本句提取主干为“上海是季节”，“冬天”是“季节”，“上海”不是“季节”，应改为“上海的冬天是美丽的季节”。

\subsubsection{修饰语和中心语搭配不当}

例：我们有一双聪明能干的手，什么造不出来？

解析：本句属于修饰语和中心语搭配不当。“聪明”不能修饰“手”，应把“聪明”改为“灵巧”。

\subsubsection{一面与两面搭配不当}

例：摄影作品拍得好坏，诗歌写得有味无味，是由一个人的思想、艺术修养决定的。

解析：句中“思想、艺术修养”的含义是确定的，意思是“有思想、艺术修养”，只能与前面的“拍得好”、“写得有味”意思搭配。这个句子要做到前后搭配，需要在“思想、艺术修养”后加“高低”一词。

\subsection{用词不当}

\begin{enumerate}[leftmargin=*]
    \item 感情色彩不当。如：他认真刻苦的学习精神，值得我们每个同学效尤。（“效尤”含贬义，可改为“学习”）
    \item 关联词使用不当。如：只要坚持下去，我们才能取得成功。（关联词语一般成对使用，“只要”一般与“就”搭配，“只有”与“才”搭配。）
    \item 意思表达不当。如：他这个人，性格直爽，为人忠厚，说起话来信口开河。（“信口开河”意为“随口乱说一气”，此处误用为“快言直语”。）
    \item 词性误用。如：从大量统计资料来看，吸烟能导致癌症是无可疑问的。（“疑问”是名词，应改为“置疑”。）
    \item 大词小用。如：每个老师日常从事的事业，是平凡而又伟大的。（“事业”应改为“工作”。）
\end{enumerate}

\subsection{语序不当}

\subsubsection{名词修饰语次序不当}

例：在新中国的建设事业中，他们发挥着无穷的蕴藏着的力量。

解析：“无穷的”应紧靠“力量”。多项定语与中心语的正确次序一般是：（1）表领属性的或时间、处所的；（2）指称或数量的短语；（3）动词或动词短语；（4）形容词或形容词短语；（5）名词或名词短语。另外，带“的”的定语放在不带“的”的定语前。

\subsubsection{动词修饰语次序不当}

例：开考半个小时后，就有人陆续交卷。

“陆续”应移到“有”的前面。多项状语次序比较复杂，须特别注意的是：（1）先时间后处所；（2）先介词结构后动词、形容词；（3）表示对象的介词结构一般紧靠中心语；（4）不要弄错修饰对象。

\subsubsection{关联词语位置不当}

例：他如果不能实事求是，事业就会受到损失。

解析：“他”应移到“如果”后面。一般来说，两个分句同一个主语时，关联词语在主语后边；不同主语时，关联词语在主语前边。

\subsection{成分残缺}

\subsubsection{缺少主语}

例：通过形式多样的活动，使广大青年学生树立了正确的人生观和价值观。

解析：“通过”与“形式多样的活动”构成介宾短语，充当状语，后一分句由“使”领起，造成该句缺少主语，“通过”或“使”保留一个即可。

\subsubsection*{缺谓语}

例：中国人民正在为建设一个现代化的社会主义强国。

“为建设一个现代化的社会主义强国”是个介词结构，只能作状语而不能作谓语。“为建设强国”怎么样？话没说完。在“强国”后面补上“而努力奋斗”，句子就有谓语了。

\subsubsection{缺少宾语中心语}

例：经过阿富汗与伊拉克战争，美国得出无须依靠大型军事基地也可以打赢一场战争。

解析：此句缺少宾语中心语，“得出”与“打赢一场战争”不能搭配，应在“战争”后面加上“的结论”。

\subsubsection{缺介词}

例：周日，嘉成老师跟他关系最好的宏涛老师去逛街了。

乍看好像没什么问题，分析一下句子成分会发现“跟他关系最好的”是修饰“宏涛老师”的定语成分，而这里想表达的是“嘉成老师和宏涛老师”去逛街了。若要修改，在“跟”字前加个“和”即可。

\subsection{成分赘余}

成分赘余，顾名思义就是句子里有多余的成分，造成了语义重复的问题。

\subsubsection{主语有多余成分}

例：全场座位座无虚席，教授的报告赢得了同学们的热烈掌声。

“座无虚席”的“座”已经是“座位”的意思了，没有必要再说一次“座位”。

\subsubsection{谓语有多余成分}

例：这位老教授的学术造诣在国内可以堪称一流。

“堪”的意思就是可以、能够，再说一个“可以”，说重了。可以改成“在国内堪称一流”或者“在国内可以称得上一流”。

\subsubsection{宾语有多余成分}

例：只有密切接触社会，联系群众，才能对国家安危和人民快乐提出具有真知灼见的意见。

“真知灼见”的“见”就是意见的意思，这里直接说“提出真知灼见”就行了，没必要加“的意见”。

\subsubsection{定语多余}

例：你去劝劝她，让她别把过去的往事放在心上。

“往事”的“往”就是“过去”的意思，“过去的往事”，显然说重复了，可以说“别把往事放在心上”，或者“别把过去的事情放在心上”。

\subsubsection{状语多余}

例：在使用电器的时候，如果一旦出现漏电现象，应该立即切断电源。

“如果”和“一旦”都是表示假设，没有必要都用，删去任意一个即可。

\subsubsection{其他多余}

例：从那以后，这个原本平静的家里，不时发生使人不安的怪事出来。

句中的“发生”就有“出来”的意思，再加“出来”就多余了，形成了补语赘余的问题，应当删掉“出来”。

\subsection{结构混乱}

\subsubsection{一个句子套用了两种结构}

例：一个人变好变坏，关键在于内因起决定作用。

“关键在于内因”与“内因起决定作用”保留一个即可。

\subsubsection{两个分句糅成一个单句}

指本来是两个分句组成的复句，却被杂糅成一个单句。

例：小张除跳舞外，兼任报幕、开场、结尾的节目还得由她编导。

“小张……兼任报幕”和“开场……由她编导”是两个分句，各有主语，其间误用顿号；“兼任”一词一直管到“节目”，造成杂糅和搭配不当的双重毛病。

\subsubsection{藕断丝连}

指一句话结构已完整，却把它的最后部分用作另一部分的开头。

例：处理好人与自然的关系，要靠政府的力量，同时也不能不发挥民间力量在舆论动员、监督检查等方面起到无可替代的作用。

“民间力量”在前部分作宾语，整个句子已经完整，但加上“在舆论……的作用”部分，又使它作了后部分的主语，整个句子的结构乱了。为突出“民间力量”，可从它后面断开，再在“在舆论……的作用”前加“民间力量”，使整个句子变成两句话。

\subsubsection{中途易辙}

一句话说了一半，忽然另起炉灶，重来一句。

例：观摩了这次关于农村经营承包合同法的庭审以后，对我们这些“村官”的法律水平有了很大的提高。

“观摩……庭审”的主语无疑是“我们”（省略了），但接下来的“对我们……的主语”显然不是“我们”，中途更换主语造成了结构混乱，删去“对”。

\subsection{不合逻辑}

\subsubsection{自相矛盾}

例：他是多少个死难者中幸免的一个。

解析：既然“幸免”，自然是没有死，怎么能说是“死难者中的一个”呢？应改为“多少人遇难了，他是幸免的一个”。

\subsubsection{并列不当}

例：他们一面拼命地往上爬，一面又不免跌落深渊。

解析：“一面……一面……”表示两件事同时进行，句中的两件事显然不是同时的，应改为“他们虽然拼命地向上爬，但是最终不免跌落深渊”。

\subsubsection{强加因果}

例：最近我这位朋友去了一趟南方回来，结果他的思想依然如故。

解析：去了南方回来思想变了，可以说是去了一趟南方的结果，现在“思想依然如故”，怎么能说是去了一趟南方的“结果”呢？

\subsubsection{主客体颠倒}

例：他从小在这长大，这里的山山水水，对他太熟悉了。

解析：“山山水水”是客观存在，是“熟悉”的客体，“他”是“熟悉”的主体。这句话正确的说法应该是“他从小在这长大，对这里的山山水水，他太熟悉了”。

\subsection{表意不明}

\subsubsection{指代不明}

如：有人主张接受，有人反对，他同意这种主张。（这种主张到底是指接受还是反对，交代不清。）

\subsubsection{介词误用}

如：奥斯特洛夫斯基的《钢铁是怎样炼成的》对于中国青年是不陌生的。（主体是人，客体是物，即作品，应是人对物不陌生。）

\section{二、常见的辨析病句的方法}

\begin{enumerate}[leftmargin=*]
    \item \textbf{紧缩法}：这是最常用的语法分析方法。先把句中的附加成分（定语、状语和补语）都去掉，紧缩出主干，检查主干是否存在成分残缺、搭配不当的语病；如果主干没问题，再检查局部，看修饰语和中心语之间的搭配有无问题，修饰语的内部是否存在语序问题。
    \item \textbf{逻辑分析法}：有的语病从语法上不好找毛病，就得从事理上进行分析，这就是逻辑意义分析法。逻辑意义分析法要从概念、判断、推理方面考虑是否得当，从句句的前后顺序、句间关系、关联词语的使用等方面考虑是否合适。例如：“该市有人不择手段伪造伪劣产品。”此句“伪造伪劣产品”是不合事理的，应改为“制造伪劣产品”或“伪造名牌产品”。辨析病句还得从语言习惯、感情色彩、语体色彩等方面仔细推敲，逐一分析，就能辨其正误，识其病因。
    \item \textbf{语感审读法}：调动语感，在审读的过程中，从感性上察觉语句的毛病。即按习惯的说法看是否别扭。如别扭则再作分析比较，辩明原因，加以修改。
    \item \textbf{造句类比法}：对句子的毛病拿捏不准时，按原句格式仿造一些浅近的、容易把握的句子加以比较，就能比较清楚地看到语病所在。
    \item \textbf{修改病句的原则}：
        \begin{itemize}
            \item 不要改变句子的原意。
            \item 修改后不能再出现新的语病。
            \item 要保持句子的简洁。
            \item 要根据实际表达需要修改，不能为了改病句而改病句，改动的字数要尽量少。切忌违背原意，另起炉灶，按自己的意志另选一个句子去代替原句。
        \end{itemize}
\end{enumerate}

\section{考场练兵}

\subsection{下列句子中，没有语病的一项是（\ ）}

A. 滥用外来语，既消解了中国文化精深而丰富的内涵，也影响了汉语表意功能的发挥。

B. 学校能否把素质教育的功夫做实，关键在于将考试建立在发展素质教育的基础上。

C. 今年，文化和旅游部继续开展“云游非遗”线上推广活动，鼓励人们走进传统文化、发现非遗之美。

D. 假如你是一粒种子，尽管是生在沃野，还是长在峭壁，都要坚强地扎根生长。


\subsection{下列病句修改不正确的一项是（\ ）}

A. 第12版《新华字典》收录了“拼购”“刷屏”等网络词语，意味着一些网络词语被更广泛地使用和接受。（将“使用和接受”改为“接受和使用”）

B. 党的十八大以来，我国推动生态环境保护之大前所未有，环境污染防治行动成效明显。（将“明显”改为“显著”）

C. 有识之士认为，能否科学设计家庭作业是实现家校良性互动、家庭作业回归育人本位的重要条件之一。（删去“否”）

D. 贯彻落实中央文化润疆精神，就要努力创作展现兵团精神的更多新时代文艺作品。（将“更多”调到“创作”之后）



\subsection{下列对病句的修改不正确的一项是（\ ）}

A. 人行道的一手扶扶，朋友间的一份鼓励，对手间的一丝宽容……这件件微薄的善举，让我内心感到由衷的温暖。（删去“由衷的”）

B. 由于采用了科学的复习方法，他的学习效率有了较大幅度的改动。（将“改动”改为“改变”）

C. 在阅读名著中，使我们明白许多做人的道理，悟出人生的真谛。（去掉“在……中”）

D. 时代的进步要求人们开阔的视野，创新的思维。（在“人们”后加“具有”）



\subsection{下列表述不正确的一项是（\ ）}

A. 文言文中的称谓语非常丰富，有自称，有对他人的爱称，敬称等。如“卿今当涂掌事”中的“卿”是古代君对臣的爱称。

B. 《卖油翁》选自《归田录》卷一，作者欧阳修，号醉翁，晚号六一居士，谥号文忠，北宋政治家、文学家，唐宋八大家之一。

C. 汉字在这些国家的应用与传播，提升了民族文化的交流，成为中国与世界各国友谊的桥梁。（该句有语病，应把“提升”改为“提高”）

D. “沿河边（跑步）”和“从昨天（开始）”中的加着重号的词都是介词，它们跟名词或代词结合在一起组成短语，表示地点、时间等。



\subsection{下列句子中没有语病、表达清楚的一项是（\ ）}

A. 环境很艰苦，条件很困难，只有我们坚定信心，顽强不屈，才能迎来人生的绿洲。

B. 王继才在孤岛上坚守32年，把最美好的年华全部献给祖国的海防事业。

C. 为了避免安全事故不再发生，学校举办了“安全伴我行”的主题演讲活动。

D. 《阿长与〈山海经〉》的结尾，表示了作者对长妈妈的深切怀念。



\subsection{阅读下面的文段，根据要求完成下面小题。}

（1）谁是我们最可爱的人呢？我们的战士，我感到他们是最可爱的人。也许还有人心里隐隐约约地说：你说的就是那些“兵”吗？他们看来是很平凡、很简单的吧，既看不出他们有什么高深的知识，又看不出他们有什么丰富的感情。可是，我要说，这是由于他跟我们的战士接触太少，还没有了解我们的战士：他们的品质是那样地纯洁和高尚，他们的意志是那样地（3）\underline{\hspace{2em}}和刚强，他们的气质是那样地淳朴和（4）\underline{\hspace{2em}}，他们的胸怀是那样地美丽和（5）\underline{\hspace{2em}}！

（2）禾稻的香气是强烈的，碾着新谷的场院辘辘地响着，多么美丽，多么（1）丰饶……没有人能够忘记她。我必定为她而战斗到底。土地，原野，我的家乡，你必须被解放！你必须站立！夜夜我看见马蹄奔驰的声音，草原的儿子在黎明的天边呼唤。这时我起来，找寻天空中北方的大熊，在它金色的光芒之下，是我的家乡。我向那边注视着，注视着，直到天边破晓。我永不能忘记，因为我答应过她，我要回到她的身边，我答应过我一定会回去。为了她，我愿付出一切。我必须看见一个更美丽的故乡出现在我的面前——或者我的坟前，而我将用我的泪水，洗去她一切的（2）wǔ huì和耻辱。

\subsubsection*{（1）请写出相应的汉字或拼音。}

（1）丰饶 \_\_\_\_\_\_ （2）wǔ huì \_\_\_\_\_\_



\subsubsection*{（2）请结合语境，将正确的词语选项填写到相应的横线上。}

A 谦逊 \quad B. 宽广 \quad C. 坚韧

（3）\_\_\_\_\_\_ （4）\_\_\_\_\_\_ （5）\_\_\_\_\_\_



\subsubsection*{（3）请选出下列表述不正确的一项。（\ ）}

A. “谁是我们最可爱的人呢？”中的“最”是副词。

B. “他们看来是很平凡、很简单的哩”中的“很”是助词。

C. “既看不出他们有什么高深的知识”中的“既”是连词。

D. “我向那边注视着”中的“向”是介词。



\subsubsection*{（4）文段中画线句子是病句，请修改后将正确的句子写在下面的横线上。}

病句：夜夜我看见马蹄奔驰的声音，草原的儿子在黎明的天边呼唤。
